\section{Chapter Summary}
\label{sec_meth_summary}

This chapter set out the problem and every method the experiments evaluate, in the order the investigation followed.
Section~\ref{sec_meth_problem} formulated grouping as the partition of $N$ typed keypoint detections into $K$ people,
scored by pose grouping accuracy (PGA) under Hungarian matching,
with a $k$NN graph over keypoint coordinates as the input representation.
Section~\ref{sec_meth_data} described the data pipeline -
Blender scenes with the person count sampled uniformly from one to five,
the COCO adapter that presents real annotations in the same form,
and the removal of the depth channel after the diagnostic of Section~\ref{sec_meth_data_depth}.
Section~\ref{sec_meth_baseline} defined the baseline that every later method must clear:
a standard GATv2 trained with a contrastive loss and read out with $k$-means.

The methods then followed the two sides of the hypothesis.
On the read-out side,
Section~\ref{sec_meth_learned_heads} introduced three learned grouping heads
and Section~\ref{sec_meth_sadmon} ten modularity-based SA-DMoN configurations,
each testing whether a stronger grouping algorithm can improve on $k$-means over the same embeddings.
On the embedding side,
Section~\ref{sec_meth_pggat} introduced PG-GAT and its three skeleton-aware modifications to the attention mechanism -
type-pair attention bias, skeleton-relative position encoding, and same-type repulsion heads -
and Sections~\ref{sec_meth_v2} to~\ref{sec_meth_trigat} described the three variants that stress that design:
visual-feature augmentation, a hyperbolic embedding space, and a triplet-based graph primitive.
Section~\ref{sec_meth_integration} closed the inference loop with the frozen HigherHRNet detector,
the $10$-pixel matching protocol,
the frozen-backbone K-head,
the autonomous pipeline,
and the detector-swap protocol for YOLO11x-pose and RTMO-L.
Section~\ref{sec_meth_evaluation} fixed the three evaluation settings,
the run-identification layer that ties every reported number to a logged run,
and the four Bayesian sweeps through which the PG-GAT configuration is selected.
Chapter~\ref{chap_4} reports the results in the same order.
